\clearpage
\section{차수, 정상폐포와 갈루아군 판정}
\label{sec:construction-galois-criteria}

\begin{learninggoal}
작도가능성의 차수 조건과 갈루아군 조건을 구분한다. 기약 사차다항식의 다섯 추이적 갈루아군 가운데 어떤 군이 작도가능한 근을 주는지 판정한다.
\end{learninggoal}

표준 시작점 $M=\{0,1\}$에서는 기준체가 $\Q$이다. 따라서 대수적 수 $\alpha$가 작도가능하려면
\[
  [\Q(\alpha):\Q]=2^r
\]
이어야 한다. 그러나 이 조건은 최소다항식이 갖는 한 근의 차수만 본다. 작도는 한 근을 얻은 뒤 그 근의 켤레들도 모두 어떤 갈루아 $2$-확장 안에 들어가야 하므로 정상폐포의 구조가 결정적이다.

\begin{theorem}[정상폐포 판정]
\label{thm:normal-closure-constructibility}
$\alpha\in\overline\Q$이고 $N$이 $\Q(\alpha)/\Q$의 정상폐포라 하자. 다음은 동치이다.
\begin{enumerate}[label=(\roman*)]
  \item $\alpha$는 자와 컴퍼스로 작도가능하다.
  \item $[N:\Q]$는 $2$의 거듭제곱이다.
  \item $\Gal(N/\Q)$는 유한 $2$-군이다.
\end{enumerate}
\end{theorem}

이는 앞 절의 따름정리를 $M=\{0,1\}$에 적용한 것이다. 따라서 실제 판정에서는 최소다항식의 차수보다 갈루아군을 계산하는 편이 더 정확하다.

\subsection{기약 사차식의 완전한 판정}

$\Q$ 위 기약 사차다항식의 갈루아군은 근 네 개에 추이적으로 작용하므로 다음 다섯 군 가운데 하나이다.
\[
  V_4,\qquad C_4,\qquad D_4,\qquad A_4,\qquad S_4.
\]

\begin{corollary}[기약 사차식과 작도가능성]
\label{cor:quartic-constructibility-groups}
$f\in\Q[X]$가 분리인 기약 사차다항식이고 $\alpha$가 한 근이라 하자. 그러면
\[
  \alpha\text{가 작도가능}
  \quad\Longleftrightarrow\quad
  \Gal(f/\Q)\in\{V_4,C_4,D_4\}.
\]
$A_4$와 $S_4$의 경우에는 어떤 근도 작도가능하지 않다.
\end{corollary}

\begin{proof}
각 군의 위수는
\[
  |V_4|=4,\quad |C_4|=4,\quad |D_4|=8,\quad
  |A_4|=12,\quad |S_4|=24.
\]
앞의 정리에 의해 분해체의 갈루아군이 $2$-군일 때 정확히 작도가능하다.
\end{proof}

\begin{center}
\renewcommand{\arraystretch}{1.2}
\begin{tabular}{cccc}
\toprule
갈루아군 & 위수 & $2$-군인가? & 근의 작도가능성 \\
\midrule
$V_4$ & $4$ & 예 & 가능 \\
$C_4$ & $4$ & 예 & 가능 \\
$D_4$ & $8$ & 예 & 가능 \\
$A_4$ & $12$ & 아니오 & 불가능 \\
$S_4$ & $24$ & 아니오 & 불가능 \\
\bottomrule
\end{tabular}
\end{center}

\begin{example}[$X^4-2$]
$\alpha=\sqrt[4]{2}$라 하자. 직접
\[
  \Q\subset\Q(\sqrt2)\subset\Q(\sqrt[4]2)
\]
라는 이차확장 탑이 있으므로 $\alpha$는 작도가능하다. 분해체는
\[
  \Q(\sqrt[4]2,i)
\]
이고 갈루아군은 $D_4$이며 위수는 $8$이다. 체확장 판정과 직접적인 반복 제곱근 구성이 서로 일치한다.
\end{example}

\subsection{차수가 사차이지만 작도할 수 없는 수}

\begin{example}
\label{ex:degree-four-nonconstructible}
다항식
\[
  f=X^4-X-1\in\Q[X]
\]
의 어떤 근도 작도가능하지 않다. 그러나 각 근의 $\Q$ 위 차수는 $4$이다.
\end{example}

\begin{proof}
법 $2$로 줄이면
\[
  \overline f=X^4+X+1\in\F_2[X].
\]
이 다항식은 $0,1$을 근으로 갖지 않는다. $\F_2[X]$의 유일한 일계수 기약 이차다항식 $X^2+X+1$의 근 $u$에 대해 $u^2=u+1$이고 $u^4=u$이므로
\[
  u^4+u+1=1\ne0.
\]
따라서 이차인수도 없고 $\overline f$는 기약이다. 그러므로 $f$는 $\Q$ 위에서 기약이며 각 근의 차수는 $4$이다.

$f$는 우울형 사차식
\[
  X^4+pX^2+qX+r
\]
에서 $p=0,q=-1,r=-1$인 경우이다. 삼차분해식은
\[
  R_f(Z)=Z^3+4Z+1.
\]
유리근 후보 $\pm1$을 대입해도 $0$이 되지 않으므로 $R_f$는 $\Q$ 위에서 기약이다. 따라서 사차 갈루아군은 $A_4$ 또는 $S_4$이다.

한편 희소 사차식의 판별식 공식에서
\[
  \Disc(X^4+bX+c)=256c^3-27b^4
\]
이므로
\[
  \Disc(f)=-256-27=-283
\]
이다. 이는 $\Q$의 제곱이 아니므로 갈루아군은 $A_4$ 안에 들어가지 않는다. 따라서
\[
  \Gal(f/\Q)\cong S_4.
\]
분해체 차수는 $24$이고 $2$의 거듭제곱이 아니므로 근들은 작도가능하지 않다.
\end{proof}

\begin{keyidea}
\[
  \deg m_\alpha=2^r
\]
는 한 근까지 가는 길의 길이만 제한한다. 작도가능성을 완전히 결정하는 것은 모든 켤레근을 함께 담는 정상폐포의 갈루아군이다.
\end{keyidea}

\subsection{작도가능체의 갈루아적 성질}

\begin{proposition}
표준 작도가능체 $\mathcal C_0$는 다음 성질을 갖는다.
\begin{enumerate}[label=(\roman*)]
  \item $\mathcal C_0/\Q$는 대수적 확장이다.
  \item 모든 $\Q$-켤레를 함께 포함하므로 $\mathcal C_0/\Q$는 정규확장이다.
  \item 표수 $0$이므로 분리확장이고, 따라서 무한 갈루아 확장이다.
  \item $\mathcal C_0$는 유한 갈루아 $2$-확장들의 합집합이다.
\end{enumerate}
\end{proposition}

\begin{proof}
모든 원소는 유한 이차확장 탑 안에 있으므로 대수적이다. $\alpha\in\mathcal C_0$이면 그 최소다항식의 분해체는 유한 갈루아 $2$-확장이고 $\mathcal C_0$ 안에 들어간다. 따라서 모든 켤레도 $\mathcal C_0$에 속한다. 표수 $0$에서 대수적 확장은 분리적이므로 정규성과 함께 갈루아 확장이 된다. 마지막 명제는 정상폐포 판정에서 직접 따른다.
\end{proof}

\begin{corollary}
$\alpha,\beta$가 작도가능한 대수적 수이면 그 정상폐포들의 합성체도 유한 갈루아 $2$-확장이다. 특히
\[
  \alpha\pm\beta,\quad \alpha\beta,\quad \alpha/\beta\ (\beta\ne0)
\]
뿐 아니라 $\alpha$와 $\beta$의 모든 켤레도 작도가능하다.
\end{corollary}

\begin{checkpoint}
\begin{enumerate}[label=(\alph*)]
  \item 기약 사차식의 삼차분해식이 기약이고 판별식이 제곱이면 갈루아군은 무엇인가? 그 근은 작도가능한가?
  \item $X^4-X-1$의 근의 차수가 $4$인데도 작도할 수 없는 핵심 이유는 무엇인가?
  \item 작도가능체가 정규확장인 이유를 정상폐포 판정으로 설명하라.
\end{enumerate}
\end{checkpoint}
